\documentclass[10pt,a4paper]{amsart}
\usepackage{amsmath,amssymb,amsthm,graphicx}
% For commutative diagrams we use Xy-pic
\usepackage[all]{xy}

% We use 2cell for 2-commutative diagrams.
\xyoption{2cell}
% Theorem environments.
%
\theoremstyle{plain}
\newtheorem{theorem}[subsection]{Theorem}
\newtheorem{proposition}[subsection]{Proposition}
\newtheorem{lemma}[subsection]{Lemma}

\theoremstyle{definition}
\newtheorem{definition}[subsection]{Definition}
\newtheorem{example}[subsection]{Example}
\newtheorem{exercise}[subsection]{Exercise}
\newtheorem{situation}[subsection]{Situation}

\theoremstyle{remark}
\newtheorem{remark}[subsection]{Remark}
\newtheorem{remarks}[subsection]{Remarks}

\numberwithin{equation}{subsection}

% Macros
%
\def\lim{\mathop{\mathrm{lim}}\nolimits}
\def\colim{\mathop{\mathrm{colim}}\nolimits}
\def\Spec{\mathop{\mathrm{Spec}}}
\def\Hom{\mathop{\mathrm{Hom}}\nolimits}
\def\Ext{\mathop{\mathrm{Ext}}\nolimits}
\def\SheafHom{\mathop{\mathcal{H}\!\mathit{om}}\nolimits}
\def\SheafExt{\mathop{\mathcal{E}\!\mathit{xt}}\nolimits}
\def\Sch{\mathit{Sch}}
\def\Mor{\mathop{\mathrm{Mor}}\nolimits}
\def\Ob{\mathop{\mathrm{Ob}}\nolimits}
\def\Sh{\mathop{\mathit{Sh}}\nolimits}
\def\NL{\mathop{N\!L}\nolimits}
\def\CH{\mathop{\mathrm{CH}}\nolimits}
\def\proetale{{pro\text{-}\acute{e}tale}}
\def\etale{{\acute{e}tale}}
\def\QCoh{\mathit{QCoh}}
\def\Ker{\mathop{\mathrm{Ker}}}
\def\Im{\mathop{\mathrm{Im}}}
\def\Coker{\mathop{\mathrm{Coker}}}
\def\Coim{\mathop{\mathrm{Coim}}}

% Boxtimes
%
\DeclareMathSymbol{\boxtimes}{\mathbin}{AMSa}{"02}

%
% Macros for moduli stacks/spaces
%
\def\QCohstack{\mathcal{QC}\!\mathit{oh}}
\def\Cohstack{\mathcal{C}\!\mathit{oh}}
\def\Spacesstack{\mathcal{S}\!\mathit{paces}}
\def\Quotfunctor{\mathrm{Quot}}
\def\Hilbfunctor{\mathrm{Hilb}}
\def\Curvesstack{\mathcal{C}\!\mathit{urves}}
\def\Polarizedstack{\mathcal{P}\!\mathit{olarized}}
\def\Complexesstack{\mathcal{C}\!\mathit{omplexes}}
% \Pic is the operator that assigns to X its picard group, usage \Pic(X)
% \Picardstack_{X/B} denotes the Picard stack of X over B
% \Picardfunctor_{X/B} denotes the Picard functor of X over B
\def\Pic{\mathop{\mathrm{Pic}}\nolimits}
\def\Picardstack{\mathcal{P}\!\mathit{ic}}
\def\Picardfunctor{\mathrm{Pic}}
\def\Deformationcategory{\mathcal{D}\!\mathit{ef}}
\newenvironment{reference}{}{}
\newenvironment{slogan}{}{}
\newenvironment{history}{}{}
\usepackage{cleveref}
\crefname{originalproposition}{original proposition}{original propositions}
\Crefname{originalproposition}{Original proposition}{Original propositions}
\makeatletter
\newlabel{proposition-polynomial}{{08PX}{source}{}{sourcepointer-1}{}}
\newlabel{proposition-polynomial@cref}{{[originalproposition][0][]08PX}{[1][0][]source}{}{}{}}
\makeatother
\crefname{originalsection}{original section}{original sections}
\Crefname{originalsection}{Original section}{Original sections}
\makeatletter
\newlabel{simplicial-section-standard}{{08N8}{source}{}{sourcepointer-2}{}}
\newlabel{simplicial-section-standard@cref}{{[originalsection][0][]08N8}{[1][0][]source}{}{}{}}
\makeatother
\crefname{originalsection}{original section}{original sections}
\Crefname{originalsection}{Original section}{Original sections}
\makeatletter
\newlabel{section-compute-L-pi-shriek}{{08PQ}{source}{}{sourcepointer-3}{}}
\newlabel{section-compute-L-pi-shriek@cref}{{[originalsection][0][]08PQ}{[1][0][]source}{}{}{}}
\makeatother
\crefname{originalsection}{original section}{original sections}
\Crefname{originalsection}{Original section}{Original sections}
\makeatletter
\newlabel{sites-cohomology-example-morphism-categories}{{08PH}{source}{}{sourcepointer-4}{}}
\newlabel{sites-cohomology-example-morphism-categories@cref}{{[originalsection][0][]08PH}{[1][0][]source}{}{}{}}
\makeatother
\crefname{originalsection}{original section}{original sections}
\Crefname{originalsection}{Original section}{Original sections}
\makeatletter
\newlabel{sites-cohomology-remark-morphism-fibred-categories}{{08PD}{source}{}{sourcepointer-5}{}}
\newlabel{sites-cohomology-remark-morphism-fibred-categories@cref}{{[originalsection][0][]08PD}{[1][0][]source}{}{}{}}
\makeatother
\crefname{originallemma}{original lemma}{original lemmas}
\Crefname{originallemma}{Original lemma}{Original lemmas}
\makeatletter
\newlabel{lemma-compute-cotangent-complex}{{08PU}{source}{}{sourcepointer-6}{}}
\newlabel{lemma-compute-cotangent-complex@cref}{{[originallemma][0][]08PU}{[1][0][]source}{}{}{}}
\makeatother
\crefname{originallemma}{original lemma}{original lemmas}
\Crefname{originallemma}{Original lemma}{Original lemmas}
\makeatletter
\newlabel{sites-cohomology-lemma-get-it-now}{{08QA}{source}{}{sourcepointer-7}{}}
\newlabel{sites-cohomology-lemma-get-it-now@cref}{{[originallemma][0][]08QA}{[1][0][]source}{}{}{}}
\makeatother
\crefname{originaldefinition}{original definition}{original definitions}
\Crefname{originaldefinition}{Original definition}{Original definitions}
\makeatletter
\newlabel{algebra-definition-naive-cotangent-complex}{{07BN}{source}{}{sourcepointer-8}{}}
\newlabel{algebra-definition-naive-cotangent-complex@cref}{{[originaldefinition][0][]07BN}{[1][0][]source}{}{}{}}
\makeatother
\crefname{originalequation}{original equation}{original equations}
\Crefname{originalequation}{Original equation}{Original equations}
\makeatletter
\newlabel{equation-comparison-map}{{08R9}{source}{}{sourcepointer-9}{}}
\newlabel{equation-comparison-map@cref}{{[originalequation][0][]08R9}{[1][0][]source}{}{}{}}
\makeatother
\crefname{originalequation}{original equation}{original equations}
\Crefname{originalequation}{Original equation}{Original equations}
\makeatletter
\newlabel{equation-triangle}{{08QS}{source}{}{sourcepointer-10}{}}
\newlabel{equation-triangle@cref}{{[originalequation][0][]08QS}{[1][0][]source}{}{}{}}
\makeatother
\crefname{originallemma}{original lemma}{original lemmas}
\Crefname{originallemma}{Original lemma}{Original lemmas}
\makeatletter
\newlabel{lemma-cotangent-complex-polynomial-algebra}{{08QH}{source}{}{sourcepointer-11}{}}
\newlabel{lemma-cotangent-complex-polynomial-algebra@cref}{{[originallemma][0][]08QH}{[1][0][]source}{}{}{}}
\makeatother
\crefname{originallemma}{original lemma}{original lemmas}
\Crefname{originallemma}{Original lemma}{Original lemmas}
\makeatletter
\newlabel{algebra-lemma-NL-polynomial-algebra}{{08Q1}{source}{}{sourcepointer-12}{}}
\newlabel{algebra-lemma-NL-polynomial-algebra@cref}{{[originallemma][0][]08Q1}{[1][0][]source}{}{}{}}
\makeatother
\crefname{originallemma}{original lemma}{original lemmas}
\Crefname{originallemma}{Original lemma}{Original lemmas}
\makeatletter
\newlabel{lemma-surjection}{{08RA}{source}{}{sourcepointer-13}{}}
\newlabel{lemma-surjection@cref}{{[originallemma][0][]08RA}{[1][0][]source}{}{}{}}
\makeatother
\crefname{originallemma}{original lemma}{original lemmas}
\Crefname{originallemma}{Original lemma}{Original lemmas}
\makeatletter
\newlabel{algebra-lemma-NL-surjection}{{07BP}{source}{}{sourcepointer-14}{}}
\newlabel{algebra-lemma-NL-surjection@cref}{{[originallemma][0][]07BP}{[1][0][]source}{}{}{}}
\makeatother
\crefname{originallemma}{original lemma}{original lemmas}
\Crefname{originallemma}{Original lemma}{Original lemmas}
\makeatletter
\newlabel{algebra-lemma-exact-sequence-NL}{{00S2}{source}{}{sourcepointer-15}{}}
\newlabel{algebra-lemma-exact-sequence-NL@cref}{{[originallemma][0][]00S2}{[1][0][]source}{}{}{}}
\makeatother
\crefname{originalsection}{original section}{original sections}
\Crefname{originalsection}{Original section}{Original sections}
\makeatletter
\newlabel{remark-make-map}{{08R7}{source}{}{sourcepointer-16}{}}
\newlabel{remark-make-map@cref}{{[originalsection][0][]08R7}{[1][0][]source}{}{}{}}
\makeatother
\crefname{originallemma}{original lemma}{original lemmas}
\Crefname{originallemma}{Original lemma}{Original lemmas}
\makeatletter
\newlabel{lemma-identify-pi-shriek}{{08PS}{source}{}{sourcepointer-17}{}}
\newlabel{lemma-identify-pi-shriek@cref}{{[originallemma][0][]08PS}{[1][0][]source}{}{}{}}
\makeatother
\crefname{originalsection}{original section}{original sections}
\Crefname{originalsection}{Original section}{Original sections}
\makeatletter
\newlabel{remark-elucidate-ss}{{08RG}{source}{}{sourcepointer-18}{}}
\newlabel{remark-elucidate-ss@cref}{{[originalsection][0][]08RG}{[1][0][]source}{}{}{}}
\makeatother
\crefname{originalsection}{original section}{original sections}
\Crefname{originalsection}{Original section}{Original sections}
\makeatletter
\newlabel{simplicial-section-complexes}{{0194}{source}{}{sourcepointer-19}{}}
\newlabel{simplicial-section-complexes@cref}{{[originalsection][0][]0194}{[1][0][]source}{}{}{}}
\makeatother
\crefname{originallemma}{original lemma}{original lemmas}
\Crefname{originallemma}{Original lemma}{Original lemmas}
\makeatletter
\newlabel{simplicial-lemma-trivial-kan-homotopy}{{08NS}{source}{}{sourcepointer-20}{}}
\newlabel{simplicial-lemma-trivial-kan-homotopy@cref}{{[originallemma][0][]08NS}{[1][0][]source}{}{}{}}
\makeatother
\crefname{originallemma}{original lemma}{original lemmas}
\Crefname{originallemma}{Original lemma}{Original lemmas}
\makeatletter
\newlabel{simplicial-lemma-homotopy-equivalence}{{08P2}{source}{}{sourcepointer-21}{}}
\newlabel{simplicial-lemma-homotopy-equivalence@cref}{{[originallemma][0][]08P2}{[1][0][]source}{}{}{}}
\makeatother
\crefname{originallemma}{original lemma}{original lemmas}
\Crefname{originallemma}{Original lemma}{Original lemmas}
\makeatletter
\newlabel{simplicial-lemma-qis-simplicial-abelian-groups}{{08P1}{source}{}{sourcepointer-22}{}}
\newlabel{simplicial-lemma-qis-simplicial-abelian-groups@cref}{{[originallemma][0][]08P1}{[1][0][]source}{}{}{}}
\makeatother
\crefname{originallemma}{original lemma}{original lemmas}
\Crefname{originallemma}{Original lemma}{Original lemmas}
\makeatletter
\newlabel{sites-cohomology-lemma-compute-by-cosimplicial-resolution}{{08Q9}{source}{}{sourcepointer-23}{}}
\newlabel{sites-cohomology-lemma-compute-by-cosimplicial-resolution@cref}{{[originallemma][0][]08Q9}{[1][0][]source}{}{}{}}
\makeatother
\crefname{originallemma}{original lemma}{original lemmas}
\Crefname{originallemma}{Original lemma}{Original lemmas}
\makeatletter
\newlabel{lemma-pi-shriek-standard}{{08PT}{source}{}{sourcepointer-24}{}}
\newlabel{lemma-pi-shriek-standard@cref}{{[originallemma][0][]08PT}{[1][0][]source}{}{}{}}
\makeatother
\title{Lichtenbaum–Schlessinger complex computes the cotangent complex through degree minus two}
\author{The Stacks Project Authors}
\date{stacks-project:cotangent.tex:local-snapshot-20261011; GNU FDL 1.2 or later}
\begin{document}
\maketitle
\noindent\textit{Editorial note.} The source mixes homological H\_2 notation with cohomological H\textasciicircum{}\{-2\}; the complex is placed in degrees -2,-1,0 as stated in2392–2417.\par\medskip
\noindent\textit{Editorial note.} Local resolution and comparison results are retained supporting work and are not counted as additional tasks.\par\medskip
\noindent\textit{Allowed standard prerequisites (editorial).}\begin{itemize}
\item Explicit prerequisite boundary. Simplicial polynomial resolutions, simplicial identities, ordinary Tor exact sequences and tensor-product spectral sequences, the cotangent transitivity framework. The explicit two-skeleton, naive comparison and degree-two local calculation are retained as authored support rather than granted free.
\end{itemize}\medskip
\ifdefined\hypertarget\hypertarget{sourcepointer-1}{}\fi
\noindent\textit{Source reference.} Original proposition 08PX: focused-seeds/AGP/sources/stacks-snapshot-20261011/tags.txt, source lines 11040--11040. Stable Stacks tag for an expressly allowed foundational or notational input; the counted construction remains the two concrete core obligations. The retained local core is not replaced by this pointer. See standard\_prerequisites.\par\medskip
\ifdefined\hypertarget\hypertarget{sourcepointer-2}{}\fi
\noindent\textit{Source reference.} Original section 08N8: focused-seeds/AGP/sources/stacks-snapshot-20261011/tags.txt, source lines 10981--10981. Stable Stacks tag for an expressly allowed foundational or notational input; the counted construction remains the two concrete core obligations. The retained local core is not replaced by this pointer. See standard\_prerequisites.\par\medskip
\ifdefined\hypertarget\hypertarget{sourcepointer-3}{}\fi
\noindent\textit{Source reference.} Original section 08PQ: focused-seeds/AGP/sources/stacks-snapshot-20261011/tags.txt, source lines 11033--11033. Stable Stacks tag for an expressly allowed foundational or notational input; the counted construction remains the two concrete core obligations. The retained local core is not replaced by this pointer. See standard\_prerequisites.\par\medskip
\ifdefined\hypertarget\hypertarget{sourcepointer-4}{}\fi
\noindent\textit{Source reference.} Original section 08PH: focused-seeds/AGP/sources/stacks-snapshot-20261011/tags.txt, source lines 11025--11025. Stable Stacks tag for an expressly allowed foundational or notational input; the counted construction remains the two concrete core obligations. The retained local core is not replaced by this pointer. See standard\_prerequisites.\par\medskip
\ifdefined\hypertarget\hypertarget{sourcepointer-5}{}\fi
\noindent\textit{Source reference.} Original section 08PD: focused-seeds/AGP/sources/stacks-snapshot-20261011/tags.txt, source lines 11021--11021. Stable Stacks tag for an expressly allowed foundational or notational input; the counted construction remains the two concrete core obligations. The retained local core is not replaced by this pointer. See standard\_prerequisites.\par\medskip
\ifdefined\hypertarget\hypertarget{sourcepointer-6}{}\fi
\noindent\textit{Source reference.} Original lemma 08PU: focused-seeds/AGP/sources/stacks-snapshot-20261011/tags.txt, source lines 11037--11037. Stable Stacks tag for an expressly allowed foundational or notational input; the counted construction remains the two concrete core obligations. The retained local core is not replaced by this pointer. See standard\_prerequisites.\par\medskip
\ifdefined\hypertarget\hypertarget{sourcepointer-7}{}\fi
\noindent\textit{Source reference.} Original lemma 08QA: focused-seeds/AGP/sources/stacks-snapshot-20261011/tags.txt, source lines 11053--11053. Stable Stacks tag for an expressly allowed foundational or notational input; the counted construction remains the two concrete core obligations. The retained local core is not replaced by this pointer. See standard\_prerequisites.\par\medskip
\ifdefined\hypertarget\hypertarget{sourcepointer-8}{}\fi
\noindent\textit{Source reference.} Original definition 07BN: focused-seeds/AGP/sources/stacks-snapshot-20261011/tags.txt, source lines 9348--9348. Stable Stacks tag for an expressly allowed foundational or notational input; the counted construction remains the two concrete core obligations. The retained local core is not replaced by this pointer. See standard\_prerequisites.\par\medskip
\ifdefined\hypertarget\hypertarget{sourcepointer-9}{}\fi
\noindent\textit{Source reference.} Original equation 08R9: focused-seeds/AGP/sources/stacks-snapshot-20261011/tags.txt, source lines 11087--11087. Stable Stacks tag for an expressly allowed foundational or notational input; the counted construction remains the two concrete core obligations. The retained local core is not replaced by this pointer. See standard\_prerequisites.\par\medskip
\ifdefined\hypertarget\hypertarget{sourcepointer-10}{}\fi
\noindent\textit{Source reference.} Original equation 08QS: focused-seeds/AGP/sources/stacks-snapshot-20261011/tags.txt, source lines 11070--11070. Stable Stacks tag for an expressly allowed foundational or notational input; the counted construction remains the two concrete core obligations. The retained local core is not replaced by this pointer. See standard\_prerequisites.\par\medskip
\ifdefined\hypertarget\hypertarget{sourcepointer-11}{}\fi
\noindent\textit{Source reference.} Original lemma 08QH: focused-seeds/AGP/sources/stacks-snapshot-20261011/tags.txt, source lines 11060--11060. Stable Stacks tag for an expressly allowed foundational or notational input; the counted construction remains the two concrete core obligations. The retained local core is not replaced by this pointer. See standard\_prerequisites.\par\medskip
\ifdefined\hypertarget\hypertarget{sourcepointer-12}{}\fi
\noindent\textit{Source reference.} Original lemma 08Q1: focused-seeds/AGP/sources/stacks-snapshot-20261011/tags.txt, source lines 11044--11044. Stable Stacks tag for an expressly allowed foundational or notational input; the counted construction remains the two concrete core obligations. The retained local core is not replaced by this pointer. See standard\_prerequisites.\par\medskip
\ifdefined\hypertarget\hypertarget{sourcepointer-13}{}\fi
\noindent\textit{Source reference.} Original lemma 08RA: focused-seeds/AGP/sources/stacks-snapshot-20261011/tags.txt, source lines 11088--11088. Stable Stacks tag for an expressly allowed foundational or notational input; the counted construction remains the two concrete core obligations. The retained local core is not replaced by this pointer. See standard\_prerequisites.\par\medskip
\ifdefined\hypertarget\hypertarget{sourcepointer-14}{}\fi
\noindent\textit{Source reference.} Original lemma 07BP: focused-seeds/AGP/sources/stacks-snapshot-20261011/tags.txt, source lines 9349--9349. Stable Stacks tag for an expressly allowed foundational or notational input; the counted construction remains the two concrete core obligations. The retained local core is not replaced by this pointer. See standard\_prerequisites.\par\medskip
\ifdefined\hypertarget\hypertarget{sourcepointer-15}{}\fi
\noindent\textit{Source reference.} Original lemma 00S2: focused-seeds/AGP/sources/stacks-snapshot-20261011/tags.txt, source lines 1024--1024. Stable Stacks tag for an expressly allowed foundational or notational input; the counted construction remains the two concrete core obligations. The retained local core is not replaced by this pointer. See standard\_prerequisites.\par\medskip
\ifdefined\hypertarget\hypertarget{sourcepointer-16}{}\fi
\noindent\textit{Source reference.} Original section 08R7: focused-seeds/AGP/sources/stacks-snapshot-20261011/tags.txt, source lines 11085--11085. Stable Stacks tag for an expressly allowed foundational or notational input; the counted construction remains the two concrete core obligations. The retained local core is not replaced by this pointer. See standard\_prerequisites.\par\medskip
\ifdefined\hypertarget\hypertarget{sourcepointer-17}{}\fi
\noindent\textit{Source reference.} Original lemma 08PS: focused-seeds/AGP/sources/stacks-snapshot-20261011/tags.txt, source lines 11035--11035. Stable Stacks tag for an expressly allowed foundational or notational input; the counted construction remains the two concrete core obligations. The retained local core is not replaced by this pointer. See standard\_prerequisites.\par\medskip
\ifdefined\hypertarget\hypertarget{sourcepointer-18}{}\fi
\noindent\textit{Source reference.} Original section 08RG: focused-seeds/AGP/sources/stacks-snapshot-20261011/tags.txt, source lines 11094--11094. Stable Stacks tag for an expressly allowed foundational or notational input; the counted construction remains the two concrete core obligations. The retained local core is not replaced by this pointer. See standard\_prerequisites.\par\medskip
\ifdefined\hypertarget\hypertarget{sourcepointer-19}{}\fi
\noindent\textit{Source reference.} Original section 0194: focused-seeds/AGP/sources/stacks-snapshot-20261011/tags.txt, source lines 1639--1639. Stable Stacks tag for an expressly allowed foundational or notational input; the counted construction remains the two concrete core obligations. The retained local core is not replaced by this pointer. See standard\_prerequisites.\par\medskip
\ifdefined\hypertarget\hypertarget{sourcepointer-20}{}\fi
\noindent\textit{Source reference.} Original lemma 08NS: focused-seeds/AGP/sources/stacks-snapshot-20261011/tags.txt, source lines 11000--11000. Stable Stacks tag for an expressly allowed foundational or notational input; the counted construction remains the two concrete core obligations. The retained local core is not replaced by this pointer. See standard\_prerequisites.\par\medskip
\ifdefined\hypertarget\hypertarget{sourcepointer-21}{}\fi
\noindent\textit{Source reference.} Original lemma 08P2: focused-seeds/AGP/sources/stacks-snapshot-20261011/tags.txt, source lines 11010--11010. Stable Stacks tag for an expressly allowed foundational or notational input; the counted construction remains the two concrete core obligations. The retained local core is not replaced by this pointer. See standard\_prerequisites.\par\medskip
\ifdefined\hypertarget\hypertarget{sourcepointer-22}{}\fi
\noindent\textit{Source reference.} Original lemma 08P1: focused-seeds/AGP/sources/stacks-snapshot-20261011/tags.txt, source lines 11009--11009. Stable Stacks tag for an expressly allowed foundational or notational input; the counted construction remains the two concrete core obligations. The retained local core is not replaced by this pointer. See standard\_prerequisites.\par\medskip
\ifdefined\hypertarget\hypertarget{sourcepointer-23}{}\fi
\noindent\textit{Source reference.} Original lemma 08Q9: focused-seeds/AGP/sources/stacks-snapshot-20261011/tags.txt, source lines 11052--11052. Stable Stacks tag for an expressly allowed foundational or notational input; the counted construction remains the two concrete core obligations. The retained local core is not replaced by this pointer. See standard\_prerequisites.\par\medskip
\ifdefined\hypertarget\hypertarget{sourcepointer-24}{}\fi
\noindent\textit{Source reference.} Original lemma 08PT: focused-seeds/AGP/sources/stacks-snapshot-20261011/tags.txt, source lines 11036--11036. Stable Stacks tag for an expressly allowed foundational or notational input; the counted construction remains the two concrete core obligations. The retained local core is not replaced by this pointer. See standard\_prerequisites.\par\medskip
\begin{lemma}
\label{lemma-compare-higher}
In the situation above denote $L$ the complex
(\ref{equation-lichtenbaum-schlessinger}).
There is a canonical map $L_{B/A} \to L$ in $D(B)$ which
induces an isomorphism $\tau_{\geq -2}L_{B/A} \to L$ in $D(B)$.
\end{lemma}

\begin{remark}[Resolutions]
\label{remark-resolution}
Let $A \to B$ be any ring map. Let us call an augmented simplicial $A$-algebra
$\epsilon : P_\bullet \to B$ a {\it resolution of $B$ over $A$} if
each $P_n$ is a polynomial algebra and $\epsilon$ is a trivial Kan fibration
of simplicial sets. If $P_\bullet \to B$ is an augmentation of a simplicial
$A$-algebra with each $P_n$ a polynomial algebra surjecting onto $B$, then
the following are equivalent
\begin{enumerate}
\item $\epsilon : P_\bullet \to B$ is a resolution of $B$ over $A$,
\item $\epsilon : P_\bullet \to B$ is a quasi-isomorphism on
associated complexes,
\item $\epsilon : P_\bullet \to B$ induces a homotopy equivalence
of simplicial sets.
\end{enumerate}
To see this use Simplicial, Lemmas
\ref{simplicial-lemma-trivial-kan-homotopy},
\ref{simplicial-lemma-homotopy-equivalence}, and
\ref{simplicial-lemma-qis-simplicial-abelian-groups}.
A resolution $P_\bullet$ of $B$ over $A$ gives a cosimplicial object
$U_\bullet$ of $\mathcal{C}_{B/A}$ as in Cohomology on Sites, Lemma
\ref{sites-cohomology-lemma-compute-by-cosimplicial-resolution}
and it follows that
$$
L\pi_!\mathcal{F} = \mathcal{F}(P_\bullet)
$$
functorially in $\mathcal{F}$, see Lemma \ref{lemma-identify-pi-shriek}.
The (formal part of the) proof of Proposition \ref{proposition-polynomial}
shows that resolutions exist. We also have seen in the first proof of
Lemma \ref{lemma-pi-shriek-standard} that the standard resolution of $B$
over $A$ is a resolution (so that this terminology doesn't lead to a conflict).
However, the argument in the proof of Proposition \ref{proposition-polynomial}
shows the existence of resolutions without appealing to the simplicial
computations in Simplicial, Section \ref{simplicial-section-standard}.
Moreover, for {\it any} choice of resolution we have a canonical isomorphism
$$
L_{B/A} = \Omega_{P_\bullet/A} \otimes_{P_\bullet, \epsilon} B
$$
in $D(B)$ by Lemma \ref{lemma-compute-cotangent-complex}. The freedom to
choose an arbitrary resolution can be quite useful.
\end{remark}

\section{Comparison with Lichtenbaum-Schlessinger}
\label{section-compare-higher}

\noindent
Let $A \to B$ be a ring map. In \cite{Lichtenbaum-Schlessinger}
there is a fairly explicit determination of $\tau_{\geq -2}L_{B/A}$
which is often used in calculations of versal deformation spaces of
singularities. The construction follows.
Choose a polynomial algebra $P$ over $A$
and a surjection $P \to B$ with kernel $I$. Choose generators
$f_t$, $t \in T$ for $I$ which induces a surjection
$F = \bigoplus_{t \in T} P \to I$ with $F$ a free $P$-module.
Let $Rel \subset F$ be the kernel of $F \to I$, in other words
$Rel$ is the set of relations among the $f_t$. Let $TrivRel \subset Rel$
be the submodule of trivial relations, i.e., the submodule of $Rel$
generated by the elements $(\ldots, f_{t'}, 0, \ldots, 0, -f_t, 0, \ldots)$.
Consider the complex of $B$-modules
\begin{equation}
\label{equation-lichtenbaum-schlessinger}
Rel/TrivRel \longrightarrow
F \otimes_P B \longrightarrow
\Omega_{P/A} \otimes_P B
\end{equation}
where the last term is placed in degree $0$. The first map is the obvious
one and the second map sends the basis element corresponding to $t \in T$
to $\text{d}f_t \otimes 1$.

\begin{definition}
\label{definition-biderivation}
Let $A \to B$ be a ring map. Let $M$ be a $(B, B)$-bimodule
over $A$. An {\it $A$-biderivation} is an $A$-linear map $\lambda : B \to M$
such that $\lambda(xy) = x\lambda(y) + \lambda(x)y$.
\end{definition}

\begin{example}
\label{example-resolution-length-two}
Let $A \to B$ be a ring map. In this example we 
will construct an ``explicit'' resolution $P_\bullet$ of $B$ over $A$ of
length $2$. To do this we follow the procedure of the proof of
Proposition \ref{proposition-polynomial}, see also the discussion in
Remark \ref{remark-resolution}.

\medskip\noindent
We choose a surjection $P_0 = A[u_i] \to B$ where $u_i$ is a set of
variables. Choose generators $f_t \in P_0$, $t \in T$ of the ideal
$\Ker(P_0 \to B)$. We choose $P_1 = A[u_i, x_t]$ with face maps
$d_0$ and $d_1$ the unique $A$-algebra maps with $d_j(u_i) = u_i$
and $d_0(x_t) = 0$ and $d_1(x_t) = f_t$. The map $s_0 : P_0 \to P_1$
is the unique $A$-algebra map with $s_0(u_i) = u_i$. It is clear that
$$
P_1 \xrightarrow{d_0 - d_1} P_0 \to B \to 0
$$
is exact, in particular the map $(d_0, d_1) : P_1 \to P_0 \times_B P_0$
is surjective. Thus, if $P_\bullet$ denotes the $1$-truncated
simplicial $A$-algebra given by $P_0$, $P_1$, $d_0$, $d_1$, and $s_0$, then
the augmentation $\text{cosk}_1(P_\bullet) \to B$ is a trivial Kan fibration.
The next step of the procedure in the proof of
Proposition \ref{proposition-polynomial}
is to choose a polynomial algebra $P_2$ and a surjection
$$
P_2 \longrightarrow \text{cosk}_1(P_\bullet)_2
$$
Recall that
$$
\text{cosk}_1(P_\bullet)_2 =
\{(g_0, g_1, g_2) \in P_1^3 \mid d_0(g_0) = d_0(g_1),
d_1(g_0) = d_0(g_2), d_1(g_1) = d_1(g_2)\}
$$
Thinking of $g_i \in P_1$ as a polynomial in $x_t$ the conditions
are
$$
g_0(0) = g_1(0),\quad
g_0(f_t) = g_2(0),\quad
g_1(f_t) = g_2(f_t)
$$
Thus $\text{cosk}_1(P_\bullet)_2$ contains the elements
$y_t = (x_t, x_t, f_t)$ and $z_t = (0, x_t, x_t)$.
Every element $G$ in $\text{cosk}_1(P_\bullet)_2$ is
of the form $G = H + (0, 0, g)$ where $H$ is in the image
of $A[u_i, y_t, z_t] \to \text{cosk}_1(P_\bullet)_2$. Here
$g \in P_1$ is a polynomial with vanishing
constant term such that $g(f_t) = 0$ in $P_0$. Observe that
\begin{enumerate}
\item $g = x_t x_{t'} - f_t x_{t'}$ and
\item $g = \sum r_t x_t$ with $r_t \in P_0$ if $\sum r_t f_t = 0$ in $P_0$
\end{enumerate}
are elements of $P_1$ of the desired form. Let
$$
Rel = \Ker(\bigoplus\nolimits_{t \in T} P_0 \longrightarrow P_0),\quad
(r_t) \longmapsto \sum r_tf_t
$$
We set $P_2 = A[u_i, y_t, z_t, v_r, w_{t, t'}]$ where
$r = (r_t) \in Rel$, with map
$$
P_2 \longrightarrow \text{cosk}_1(P_\bullet)_2
$$
given by $y_t  \mapsto (x_t, x_t, f_t)$,
$z_t \mapsto (0, x_t, x_t)$,
$v_r \mapsto (0, 0, \sum r_t x_t)$, and
$w_{t, t'} \mapsto (0, 0, x_t x_{t'} - f_t x_{t'})$. A calculation
(omitted) shows that this map is surjective. Our choice of the
map displayed above determines the maps $d_0, d_1, d_2 : P_2 \to P_1$.
Finally, the procedure in the proof of
Proposition \ref{proposition-polynomial}
tells us to choose the maps $s_0, s_1 : P_1 \to P_2$ lifting the two
maps $P_1 \to \text{cosk}_1(P_\bullet)_2$. It is clear that we can take
$s_i$ to be the unique $A$-algebra maps determined by
$s_0(x_t) = y_t$ and $s_1(x_t) = z_t$.
\end{example}








\section{Functoriality}
\label{section-functoriality}

\noindent
In this section we consider a commutative square
\begin{equation}
\label{equation-commutative-square}
\vcenter{
\xymatrix{
B \ar[r] & B' \\
A \ar[u] \ar[r] & A' \ar[u]
}
}
\end{equation}
of ring maps. We claim there is a canonical $B$-linear map of complexes
$$
L_{B/A} \longrightarrow L_{B'/A'}
$$
associated to this diagram. Namely, if $P_\bullet \to B$ is the
standard resolution of $B$ over $A$ and $P'_\bullet \to B'$ is the
standard resolution of $B'$ over $A'$, then there is a canonical map
$P_\bullet \to P'_\bullet$
of simplicial $A$-algebras compatible with the augmentations
$P_\bullet \to B$ and $P'_\bullet \to B'$. This can be seen in terms
of the construction of standard resolutions in
Simplicial, Section \ref{simplicial-section-standard}
but in the special case at hand it probably suffices to say simply
that the maps
$$
P_0 = A[B] \longrightarrow A'[B'] = P'_0,\quad
P_1 = A[A[B]] \longrightarrow A'[A'[B']] = P'_1,
$$
and so on are given by the given maps $A \to A'$ and $B \to B'$.
The desired map $L_{B/A} \to L_{B'/A'}$ then comes from the associated
maps $\Omega_{P_n/A} \to \Omega_{P'_n/A'}$.

\medskip\noindent
Another description of the functoriality map can be given as follows.
Let $\mathcal{C} = \mathcal{C}_{B/A}$ and $\mathcal{C}' = \mathcal{C}_{B'/A}'$
be the categories considered in Section \ref{section-compute-L-pi-shriek}.
There is a functor
$$
u : \mathcal{C} \longrightarrow \mathcal{C}',\quad
(P, \alpha) \longmapsto (P \otimes_A A', c \circ (\alpha \otimes 1))
$$
where $c : B \otimes_A A' \to B'$ is the obvious map. As discussed in
Cohomology on Sites, Example
\ref{sites-cohomology-example-morphism-categories}
we obtain a morphism of topoi $g : \Sh(\mathcal{C}) \to \Sh(\mathcal{C}')$
and a commutative diagram of maps of ringed topoi
\begin{equation}
\label{equation-double-square}
\vcenter{
\xymatrix{
(\Sh(\mathcal{C}'), \underline{B}) \ar[d]_{\pi'} &
(\Sh(\mathcal{C}'), \underline{B'}) \ar[d]_{\pi'} \ar[l]^h &
(\Sh(\mathcal{C}), \underline{B'}) \ar[d]_\pi \ar[l]^g \\
(\Sh(*), B) &
(\Sh(*), B') \ar[l]_f &
(\Sh(*), B') \ar[l]
}
}
\end{equation}
Here $h$ is the identity on underlying topoi and given by the ring map
$B \to B'$ on sheaves of rings. 
By Cohomology on Sites, Remark
\ref{sites-cohomology-remark-morphism-fibred-categories}
given $\mathcal{F}$ on $\mathcal{C}$ and $\mathcal{F}'$ on $\mathcal{C}'$
and a transformation $t : \mathcal{F} \to g^{-1}\mathcal{F}'$
we obtain a canonical map $L\pi_!(\mathcal{F}) \to L\pi'_!(\mathcal{F}')$.
If we apply this to the sheaves
$$
\mathcal{F} : (P, \alpha) \mapsto \Omega_{P/A} \otimes_P B,\quad
\mathcal{F}' : (P', \alpha') \mapsto \Omega_{P'/A'} \otimes_{P'} B',
$$
and the transformation $t$ given by the canonical maps
$$
\Omega_{P/A} \otimes_P B \longrightarrow
\Omega_{P \otimes_A A'/A'} \otimes_{P \otimes_A A'} B'
$$
to get a canonical map
$$
L\pi_!(\Omega_{\mathcal{O}/A} \otimes_\mathcal{O} \underline{B})
\longrightarrow
L\pi'_!(\Omega_{\mathcal{O}'/A'} \otimes_{\mathcal{O}'} \underline{B'})
$$
By Lemma \ref{lemma-compute-cotangent-complex} this gives
$L_{B/A} \to L_{B'/A'}$. We omit the verification that this map
agrees with the map defined above in terms of simplicial
resolutions.

\begin{lemma}
\label{lemma-flat-base-change}
Assume (\ref{equation-commutative-square}) induces a quasi-isomorphism
$B \otimes_A^\mathbf{L} A' = B'$. Then, with notation as in
(\ref{equation-double-square}) and
$\mathcal{F}' \in \textit{Ab}(\mathcal{C}')$,
we have $L\pi_!(g^{-1}\mathcal{F}') = L\pi'_!(\mathcal{F}')$.
\end{lemma}

\begin{proof}
We use the results of Remark \ref{remark-resolution} without
further mention. We will apply Cohomology on Sites, Lemma
\ref{sites-cohomology-lemma-get-it-now}. Let $P_\bullet \to B$ be a resolution.
If we can show that $u(P_\bullet) = P_\bullet \otimes_A A' \to B'$
is a quasi-isomorphism, then we are done. The complex of $A$-modules
$s(P_\bullet)$ associated to $P_\bullet$ (viewed as a simplicial $A$-module)
is a free $A$-module resolution of $B$. Namely, $P_n$ is a free $A$-module and
$s(P_\bullet) \to B$ is a quasi-isomorphism. Thus $B \otimes_A^\mathbf{L} A'$
is computed by $s(P_\bullet) \otimes_A A' = s(P_\bullet \otimes_A A')$.
Therefore the assumption of the lemma signifies that
$\epsilon' : P_\bullet \otimes_A A' \to B'$ is a quasi-isomorphism.
\end{proof}

\begin{lemma}
\label{lemma-relation-with-naive-cotangent-complex}
Let $A \to B$ be a ring map. Then $\tau_{\geq -1}L_{B/A}$
is canonically quasi-isomorphic to the naive cotangent complex.
\end{lemma}

\begin{proof}
Consider $P = A[B] \to B$ with kernel $I$. The naive cotangent
complex $\NL_{B/A}$ of $B$ over $A$ is the complex
$I/I^2 \to \Omega_{P/A} \otimes_P B$,
see Algebra, Definition \ref{algebra-definition-naive-cotangent-complex}.
Observe that in (\ref{equation-comparison-map}) we have already
constructed a canonical map
$$
c : \NL_{B/A} \longrightarrow \tau_{\geq -1}L_{B/A}
$$
Consider the distinguished triangle (\ref{equation-triangle})
$$
L_{P/A} \otimes_P^\mathbf{L} B \to L_{B/A} \to L_{B/P} \to 
(L_{P/A} \otimes_P^\mathbf{L} B)[1]
$$
associated to the ring maps $A \to A[B] \to B$. We know that
$L_{P/A} = \Omega_{P/A}[0] = \NL_{P/A}$ in $D(P)$
(Lemma \ref{lemma-cotangent-complex-polynomial-algebra}
and
Algebra, Lemma \ref{algebra-lemma-NL-polynomial-algebra})
and that
$\tau_{\geq -1}L_{B/P} = I/I^2[1] = \NL_{B/P}$ in $D(B)$
(Lemma \ref{lemma-surjection} and
Algebra, Lemma \ref{algebra-lemma-NL-surjection}).
To show $c$ is a quasi-isomorphism it suffices by
Algebra, Lemma \ref{algebra-lemma-exact-sequence-NL}
and the long exact cohomology sequence associated to the
distinguished triangle
to show that the maps $L_{P/A} \to L_{B/A} \to L_{B/P}$ are compatible
on cohomology groups with the corresponding maps
$\NL_{P/A} \to \NL_{B/A} \to \NL_{B/P}$
of the naive cotangent complex. We omit the verification.
\end{proof}

\begin{remark}
\label{remark-explicit-comparison-map}
We can make the comparison map of
Lemma \ref{lemma-relation-with-naive-cotangent-complex}
explicit in the following way.
Let $P_\bullet$ be the standard resolution of $B$
over $A$.
Let $I = \Ker(A[B] \to B)$.
Recall that $P_0 = A[B]$. The map of the
lemma is given by the commutative diagram
$$
\xymatrix{
L_{B/A} \ar[d] & \ldots \ar[r] &
\Omega_{P_2/A} \otimes_{P_2} B
\ar[r] \ar[d] &
\Omega_{P_1/A} \otimes_{P_1} B
\ar[r] \ar[d] &
\Omega_{P_0/A} \otimes_{P_0} B
\ar[d] \\
\NL_{B/A} & \ldots \ar[r] &
0 \ar[r] & 
I/I^2 \ar[r] &
\Omega_{P_0/A} \otimes_{P_0} B
}
$$
We construct the downward arrow with target $I/I^2$
by sending $\text{d}f \otimes b$ to the class of
$(d_0(f) - d_1(f))b$ in $I/I^2$. Here $d_i : P_1 \to P_0$,
$i = 0, 1$ are the two face maps of the simplicial structure.
This makes sense as $d_0 - d_1$ maps $P_1$ into $I = \Ker(P_0 \to B)$.
We omit the verification that this rule is well defined.
Our map is compatible with the differential
$\Omega_{P_1/A} \otimes_{P_1} B \to \Omega_{P_0/A} \otimes_{P_0} B$
as this differential maps $\text{d}f \otimes b$ to
$\text{d}(d_0(f) - d_1(f)) \otimes b$. Moreover, the differential
$\Omega_{P_2/A} \otimes_{P_2} B \to \Omega_{P_1/A} \otimes_{P_1} B$
maps $\text{d}f \otimes b$ to $\text{d}(d_0(f) - d_1(f) + d_2(f)) \otimes b$
which are annihilated by our downward arrow. Hence a map of complexes.
We omit the verification that this is the same as the map of
Lemma \ref{lemma-relation-with-naive-cotangent-complex}.
\end{remark}

\begin{remark}
\label{remark-surjection}
Adopt notation as in Remark \ref{remark-make-map}. The arguments given
there show that the differential
$$
H_2(\mathcal{C}_{B/A}, \mathcal{J}/\mathcal{J}^2)
\longrightarrow
H_0(\mathcal{C}_{B/A}, \text{Tor}_1^\mathcal{O}(\mathcal{J}, \underline{B}))
$$
of the spectral sequence is an isomorphism. Let $\mathcal{C}'_{B/A}$
denote the full subcategory of $\mathcal{C}_{B/A}$ consisting of surjective
maps $P \to B$. The agreement of the cotangent complex with the naive
cotangent complex (Lemma \ref{lemma-relation-with-naive-cotangent-complex})
shows that we have an exact sequence of sheaves
$$
0 \to \underline{H_1(L_{B/A})} \to
\mathcal{J}/\mathcal{J}^2 \xrightarrow{\text{d}} \Omega \to
\underline{H_2(L_{B/A})} \to 0
$$
on $\mathcal{C}'_{B/A}$. It follows that $\Ker(d)$ and
$\Coker(d)$ on the whole category $\mathcal{C}_{B/A}$ have
vanishing higher homology groups, since
these are computed by the homology groups of constant simplicial abelian
groups by Lemma \ref{lemma-identify-pi-shriek}. Hence we conclude
that
$$
H_n(\mathcal{C}_{B/A}, \mathcal{J}/\mathcal{J}^2) \to H_n(L_{B/A})
$$
is an isomorphism for all $n \geq 2$. Combined with the remark above
we obtain the formula
$H_2(L_{B/A}) =
H_0(\mathcal{C}_{B/A}, \text{Tor}_1^\mathcal{O}(\mathcal{J}, \underline{B}))$.
\end{remark}

\begin{remark}
\label{remark-elucidate-degree-two}
Let $A \to B$ be a ring map. Let $P_\bullet$ be a resolution of
$B$ over $A$ (Remark \ref{remark-resolution}).
Set $J_n = \Ker(P_n \to B)$. Note that
$$
\text{Tor}_2^{P_n}(B, B) = 
\text{Tor}_1^{P_n}(J_n, B) =
\Ker(J_n \otimes_{P_n} J_n \to J_n^2).
$$
Hence $H_2(L_{B/A})$ is canonically equal to
$$
\Coker(\text{Tor}_2^{P_1}(B, B) \to \text{Tor}_2^{P_0}(B, B))
$$
by Remark \ref{remark-surjection}. To make this more explicit we choose
$P_2$, $P_1$, $P_0$ as in Example \ref{example-resolution-length-two}.
We claim that
$$
\text{Tor}_2^{P_1}(B, B) =
\wedge^2(\bigoplus\nolimits_{t \in T} B)\ \oplus
\ \bigoplus\nolimits_{t \in T} J_0\ \oplus
\ \text{Tor}_2^{P_0}(B, B)
$$
Namely, the basis elements $x_t \wedge x_{t'}$ of the first summand
corresponds to the element $x_t \otimes x_{t'} - x_{t'} \otimes x_t$
of $J_1 \otimes_{P_1} J_1$. For $f \in J_0$ the element $x_t \otimes f$
of the second summand corresponds to the element
$x_t \otimes s_0(f) - s_0(f) \otimes x_t$ of $J_1 \otimes_{P_1} J_1$.
Finally, the map $\text{Tor}_2^{P_0}(B, B) \to \text{Tor}_2^{P_1}(B, B)$
is given by $s_0$. The map
$d_0 - d_1 : \text{Tor}_2^{P_1}(B, B) \to \text{Tor}_2^{P_0}(B, B)$
is zero on the last summand, maps $x_t \otimes f$ to
$f \otimes f_t - f_t \otimes f$, and maps $x_t \wedge x_{t'}$
to $f_t \otimes f_{t'} - f_{t'} \otimes f_t$. All in all we conclude
that there is an exact sequence
$$
\wedge^2_B(J_0/J_0^2) \to \text{Tor}_2^{P_0}(B, B) \to H^{-2}(L_{B/A}) \to 0
$$
In this way we obtain a direct proof of a consequence of Quillen's spectral
sequence discussed in Remark \ref{remark-elucidate-ss}.
\end{remark}

\begin{lemma}
\label{lemma-polynomial-ring-unique}
Let $P = A[S]$ be a polynomial ring over $A$. Let $M$ be a $(P, P)$-bimodule
over $A$. Given $m_s \in M$ for $s \in S$, there exists a unique
$A$-biderivation $\lambda : P \to M$ mapping $s$ to $m_s$ for $s \in S$.
\end{lemma}

\begin{proof}
We set
$$
\lambda(s_1 \ldots s_t) =
\sum s_1 \ldots s_{i - 1} m_{s_i} s_{i + 1} \ldots s_t
$$
in $M$. Extending by $A$-linearity we obtain a biderivation.
\end{proof}

\begin{proof}
Let $P_\bullet \to B$ be a resolution of $B$ over $A$
(Remark \ref{remark-resolution}). We will identify $L_{B/A}$ with
$\Omega_{P_\bullet/A} \otimes B$. To construct the map we
make some choices.

\medskip\noindent
Choose an $A$-algebra map $\psi : P_0 \to P$ compatible with the
given maps $P_0 \to B$ and $P \to B$.

\medskip\noindent
Write $P_1 = A[S]$ for some set $S$. For $s \in S$ we may write
$$
\psi(d_0(s) - d_1(s)) = \sum p_{s, t} f_t
$$
for some $p_{s, t} \in P$. Think of $F = \bigoplus_{t \in T} P$
as a $(P_1, P_1)$-bimodule via the maps $(\psi \circ d_0, \psi \circ d_1)$.
By Lemma \ref{lemma-polynomial-ring-unique} we obtain a unique
$A$-biderivation $\lambda : P_1 \to F$ mapping $s$ to the vector with
coordinates $p_{s, t}$. By construction the composition
$$
P_1 \longrightarrow F \longrightarrow P
$$
sends $f \in P_1$ to $\psi(d_0(f) - d_1(f))$ because the map
$f \mapsto \psi(d_0(f) - d_1(f))$ is an $A$-biderivation agreeing with
the composition on generators.

\medskip\noindent
For $g \in P_2$ we claim that $\lambda(d_0(g) - d_1(g) + d_2(g))$
is an element of $Rel$. Namely, by the last remark of the previous
paragraph the image of $\lambda(d_0(g) - d_1(g) + d_2(g))$ in $P$ is
$$
\psi((d_0 - d_1)(d_0(g) - d_1(g) + d_2(g)))
$$
which is zero by Simplicial, Section \ref{simplicial-section-complexes}).

\medskip\noindent
The choice of $\psi$ determines a map
$$
\text{d}\psi \otimes 1 :
\Omega_{P_0/A} \otimes B
\longrightarrow
\Omega_{P/A} \otimes B
$$
Composing $\lambda$ with the map $F \to F \otimes B$ gives a
usual $A$-derivation as the two $P_1$-module structures on
$F \otimes B$ agree. Thus $\lambda$ determines a map
$$
\overline{\lambda} :
\Omega_{P_1/A} \otimes B
\longrightarrow
F \otimes B
$$
Finally, We obtain a $B$-linear map
$$
q :
\Omega_{P_2/A} \otimes B
\longrightarrow
Rel/TrivRel
$$
by mapping $\text{d}g$ to the class of $\lambda(d_0(g) - d_1(g) + d_2(g))$
in the quotient.

\medskip\noindent
The diagram
$$
\xymatrix{
\Omega_{P_3/A} \otimes B \ar[r] \ar[d] &
\Omega_{P_2/A} \otimes B \ar[r] \ar[d]_q &
\Omega_{P_1/A} \otimes B \ar[r] \ar[d]_{\overline{\lambda}} &
\Omega_{P_0/A} \otimes B \ar[d]_{\text{d}\psi \otimes 1} \\
0 \ar[r] &
Rel/TrivRel \ar[r] &
F \otimes B \ar[r] &
\Omega_{P/A} \otimes B
}
$$
commutes (calculation omitted) and we obtain the map of the lemma.
By Remark \ref{remark-explicit-comparison-map} and
Lemma \ref{lemma-relation-with-naive-cotangent-complex} we see that this map
induces isomorphisms $H_1(L_{B/A}) \to H_1(L)$ and $H_0(L_{B/A}) \to H_0(L)$.

\medskip\noindent
It remains to see that our map $L_{B/A} \to L$ induces an isomorphism
$H_2(L_{B/A}) \to H_2(L)$. Choose a resolution of $B$ over $A$ with
$P_0 = P = A[u_i]$ and then $P_1$ and $P_2$ as in
Example \ref{example-resolution-length-two}.
In Remark \ref{remark-elucidate-degree-two} we have constructed an exact
sequence
$$
\wedge^2_B(J_0/J_0^2) \to \text{Tor}_2^{P_0}(B, B) \to H^{-2}(L_{B/A}) \to 0
$$
where $P_0 = P$ and $J_0 = \Ker(P \to B) = I$.
Calculating the Tor group using the short exact sequences
$0 \to I \to P \to B \to 0$ and $0 \to Rel \to F \to I \to 0$
we find that
$\text{Tor}_2^P(B, B) = \Ker(Rel \otimes B \to F \otimes B)$.
The image of the map $\wedge^2_B(I/I^2) \to \text{Tor}_2^P(B, B)$
under this identification is exactly the image of $TrivRel \otimes B$.
Thus we see that $H_2(L_{B/A}) \cong H_2(L)$.

\medskip\noindent
Finally, we have to check that our map $L_{B/A} \to L$ actually induces
this isomorphism. We will use the notation and results discussed in
Example \ref{example-resolution-length-two} and
Remarks \ref{remark-elucidate-degree-two} and \ref{remark-surjection}
without further mention. Pick an element $\xi$ of
$\text{Tor}_2^{P_0}(B, B) = \Ker(I \otimes_P I \to I^2)$.
Write $\xi = \sum h_{t', t}f_{t'} \otimes f_t$ for some
$h_{t', t} \in P$. Tracing through the exact sequences above we
find that $\xi$ corresponds to the image in $Rel \otimes B$
of the element $r \in Rel \subset F = \bigoplus_{t \in T} P$ with
$t$th coordinate $r_t = \sum_{t' \in T} h_{t', t}f_{t'}$.
On the other hand, $\xi$ corresponds to the element of
$H_2(L_{B/A}) = H_2(\Omega)$ which is the image
via $\text{d} : H_2(\mathcal{J}/\mathcal{J}^2) \to H_2(\Omega)$
of the boundary of $\xi$ under the $2$-extension
$$
0 \to
\text{Tor}_2^\mathcal{O}(\underline{B}, \underline{B})
\to 
\mathcal{J} \otimes_\mathcal{O} \mathcal{J} \to \mathcal{J}
\to
\mathcal{J}/\mathcal{J}^2 \to 0
$$
We compute the successive transgressions of our element. First we have
$$
\xi = (d_0 - d_1)(- \sum s_0(h_{t', t} f_{t'}) \otimes x_t)
$$
and next we have
$$
\sum s_0(h_{t', t} f_{t'}) x_t = d_0(v_r) - d_1(v_r) + d_2(v_r)
$$
by our choice of the variables $v$ in
Example \ref{example-resolution-length-two}.
We may choose our map $\lambda$ above such that
$\lambda(u_i) = 0$ and $\lambda(x_t) = - e_t$ where $e_t \in F$
denotes the basis vector corresponding to $t \in T$.
Hence the construction of our map $q$ above sends $\text{d}v_r$ to
$$
\lambda(\sum s_0(h_{t', t} f_{t'}) x_t) =
\sum\nolimits_t \left(\sum\nolimits_{t'} h_{t', t}f_{t'}\right) e_t
$$
matching the image of $\xi$ in $Rel \otimes B$ (the two minus signs
we found above cancel out). This agreement finishes the proof.
\end{proof}

\bibliographystyle{plain}
\bibliography{my}
\end{document}
